\scanpage{83}
% letterhead

U N I V E R S I T E\\ DES SCIENCES ET TECHNIQUES\\ \uline{DU LANGUEDOC}\\
Montpellier, le 1er décembre 1977\\

JP/MD\\ Diffusion : Liste A\\ n\textsuperscript{o} 6826

AVIS DE SOUTENANCE D'UNE THESE DE DOCTORAT DE SPECIALITE\\ \uline{E C O L O G I E}
% printed form

\uline{Nom et prénom du candidat} : PONT Didier

\uline{Titre de la thèse} : Recherches sur l'évolution saisonnière du peuplement de Copépodes, Cladocènes et Ostracodes des rizières de Camargue.

\uline{Date, heure et lieu de soutenance prévus} : 9 décembre 1977 à 9 h 30, salle de cours n\textsuperscript{o} 10

\uline{Nom du Directeur de thèse} : Monsieur A. CHAMPEAU

\uline{Désignation du laboratoire où la thèse est préparée} : Centre d'écologie de Camargue.

\uline{Analyse de l'ouvrage} : Après avoir mis en évidence les grandes lignes de la structure et de l'évolution saisonnière du peuplement de crustacés des rizières camarguaises sur un grand nombre de stations, l'auteur analyse les résultats obtenus sur 3 stations suivies régulièrement et met en évidence les facteurs de milieu qui conditionnent le fonctionnement du peuplement zooplanctonique.

\uline{Durée approximative des recherches ayant abouti à la thèse} : 2 ans

\uline{Composition du jury proposé} : M. C. SAUVAGE - Président\\ M. AMANIEU\\ M. CHAMPEAU\\ M. DECAMPS\\ M. HEURTEAUX

Le Secrétaire Général,\\
\handwritten{\unsure{E Siau|E Si}}\\
\uline{E. SIAU}

MR LE PROFESSEUR GROTHENDIECK A.\\ MATHEMATIQUES \hfill 5
% archive sheet number: 82
